%!TEX root = ../hutbthesis_main.tex

\chapter{总结与展望}

\section{总结}

本文围绕机场飞机牵引与自动对接作业需求，以自动牵引车和飞机协同控制为研究对象，基于ROS与Gazebo联合仿真平台，开展了作业问题抽象、车辆运动学建模、感知与控制方法设计以及仿真实现等方面的研究。

首先，本文对机场牵引作业场景进行了分析，将飞机推出与前起落架对接过程抽象为低速、高精度、多约束的车辆控制问题。在此基础上，建立了双桥转向牵引车平面运动学模型，分析前桥单独转向、前后桥反向转向和同向转向等典型模式，并推导速度、转角和曲率之间的约束关系，为后续自动接近与对接控制提供数学基础。

其次，本文将飞机辅助感知与自动接近控制合并为一条完整技术链路。通过在飞机前起落架区域设置可识别标识，牵引车能够获得目标身份、相对距离、横向偏差和姿态偏差；在此基础上，结合近距测量和安全阈值判断，建立目标确认、误差提取、安全判定和终端停车条件。

再次，本文设计了基于相对位姿误差的分层控制方法。对接过程被划分为搜索等待、自动接近、姿态对准、微动对接和停车保持等阶段，并采用“模式管理---局部规划---误差反馈---执行器映射”的结构，使牵引车能够从目标识别开始，逐步完成低速接近、方向修正和最终停车。

最后，本文完成了ROS-Gazebo联合仿真系统的搭建。通过SolidWorks车辆模型、URDF结构描述、Gazebo场景、控制器配置和ROS节点通信，实现牵引车模型加载、转向控制、轮速控制、传感器挂载和自动接近对接流程验证。仿真结果表明，该系统能够支持牵引车在机场作业场景中的低速运动和对接过程分析，为后续实物样机开发、控制参数优化和传感器融合研究提供参考。

总体来看，本文完成了自动牵引车与飞机对接任务从理论建模、感知控制方法到仿真实现的初步研究，验证了基于ROS与Gazebo开展机场牵引车联合仿真的可行性。

\section{展望}

虽然本文完成了自动牵引对接系统的基本设计与仿真验证，但受研究周期和实验条件限制，仍存在一些需要进一步完善的方面。

第一，建模/车辆设计阶段耗时过长且采用了非常高精度的建模（首轮导出.stl模型文件达到了1.8G），直接导致模型预览困难、仿真模拟效率低下的问题；同时，本文对车辆动力学、轮胎接触、牵引连接机构和飞机受力过程进行了简化，后续需要建立更加完整的牵引车---飞机组合体动力学模型，使仿真结果更接近真实机场作业工况。

第二，受设备性能限制，没有考虑多变气候环境对车辆飞行器之间的影响，后续可进一步完善感知系统。可引入更复杂的视觉识别、点云处理和多传感器融合方法，提高系统在光照变化、遮挡和复杂环境下的可靠性。

第三，此次模拟环境较为简单，只放置了一架飞行器和牵引车模型，对现今大型国际机场起到的参考作用很小，后续可结合模型预测控制、动态避障算法和智能优化方法，在添加多机协作的同时，提高系统在动态障碍物、多车辆协同和复杂机场场景下的适应能力\textsuperscript{[20,25,28]}。

综上，机场自动牵引与飞机对接是一个融合车辆控制、机器人感知、路径规划和系统仿真的综合性课题。本文完成了基础平台和核心方法的搭建，后续仍可在模型精细化、感知鲁棒性、控制智能化和工程实用化方面继续深入研究。
